\begin{figure}[htbp]
    \centering
    \resizebox{0.92\textwidth}{!}{%
    \begin{tikzpicture}[
        node distance=1.5cm,
        font=\sffamily\small,
        box/.style={rectangle, draw=black, thick, minimum width=4cm, minimum height=0.8cm, align=center},
        kernel/.style={box, fill=orange!10},
        user/.style={box, fill=blue!10},
        hw/.style={box, fill=gray!20},
        arrow/.style={->, thick, >=stealth},
        label/.style={font=\scriptsize\itshape, text=red!80!black}
    ]

    % Nodes
    \node[hw] (nic) {Network Interface Card (NIC)};
    \node[kernel, below=of nic] (dma) {DMA to Ring Buffer};
    \node[kernel, below=of dma] (irq) {Hard IRQ \& NAPI Poll};
    \node[kernel, below=of irq] (stack) {\texttt{sk\_buff} Allocation \& TCP/IP Stack};
    \node[kernel, below=of stack] (socket) {Socket Receive Queue};
    \node[user,   below=of socket] (app) {User Application (\texttt{recvfrom})};

    % Arrows
    \draw[arrow] (nic) -- (dma);
    \draw[arrow] (dma) -- node[right, near start, xshift=0.2cm] {Interrupt} (irq);
    \draw[arrow] (irq) -- (stack);
    \draw[arrow] (stack) -- (socket);
    \draw[arrow] (socket) -- node[left, pos=0.5, xshift=-0.2cm, black] {System Call Buffer Copy} (app);

    % Jitter annotations
    \draw[<-, color=red] (dma.east) -- +(1,0) node[right, label] {Interrupt Latency};
    \draw[<-, color=red] (stack.east) -- +(1,0) node[right, label] {Memory Allocation Overhead};
    \draw[<-, color=red] (socket.south east) to[out=-20, in=20] node[right, label, pos=0.5] {Context Switch Jitter} (app.north east);

    % Braces
    \draw [decorate, decoration={brace, amplitude=10pt, mirror}, xshift=-0.5cm]
    (dma.north west) -- (socket.south west) node [black, midway, xshift=-1.5cm, rotate=90, align=center] {Linux Kernel\\Space};

    \end{tikzpicture}%
    }
    \caption{The standard Linux kernel network ingress pipeline. A packet must traverse multiple layers of abstraction—including DMA transfers, hardware interrupts, and TCP/IP stack logic—before reaching the application. Each stage, particularly the transition from kernel to user space via system calls, introduces non-deterministic delays that manifest as microsecond-scale jitter in HFT environments.}
    \label{fig:kernel_network_pipeline}
\end{figure}
